% Generated by scripts/assemble.py; corrections belong in source/corrections.json.
\examhead{26}{中国矿业大学(北京) 2025 年数学分析试卷}{30}
% p031-b01; source PDF page 31
\par\Needspace{5\baselineskip}
\blockanchor{p031-b01}
\majorquestion{一}填空题. 每题 5 分, 共 40 分.\par
% p031-b02; source PDF page 31
\blockanchor{p031-b02}
\question{1}求极限  $\lim_{n\to\infty}\frac{\int_{1}^{n}\ln\left(1+\frac{1}{\sqrt{x}}\right)\mathrm{d}x}{\sqrt{n}}=$  \_\_\_\_.\par
% p031-b03; source PDF page 31
\blockanchor{p031-b03}
\question{2}曲面 $x^{2} + yz + z^{2} - 1 = 0$ 在点 $(1,1,0)$ 处的切平面为 \_\_\_\_.\par
% p031-b04; source PDF page 31
\blockanchor{p031-b04}
\question{3}求不定积分  $\int(x\ln x)^{\frac{3}{2}}(\ln x+1)\mathrm{d}x=$  \_\_\_\_.\par
% p031-b05; source PDF page 31
\blockanchor{p031-b05}
\question{4}求 $\sum_{n=1}^{\infty} \frac{1}{n \cdot 2^n} = \underline{\quad}$.\par
% p031-b06; source PDF page 31
\blockanchor{p031-b06}
\question{5}已知 $L$ 为 $x^{2} + y^{2} = 2$ 的第一象限的部分取逆时针, 求 $\int_{L} x \mathrm{~d}y - 2y \mathrm{~d}x =$ \_\_\_\_.\par
% p031-b07; source PDF page 31
\blockanchor{p031-b07}
\question{6}$x = \frac{3}{4}$ 是Riemann函数的\_\_\_\_点(填连续, 可去间断, 跳跃间断或第二类间断).\par
% p031-b08; source PDF page 31
\blockanchor{p031-b08}
\question{7}计算 $\lim_{x\to \infty}\left(\cos \frac{1}{x} +\sin \frac{2}{x}\right)^x =$ \_\_\_\_.\par
% p031-b09; source PDF page 31
\blockanchor{p031-b09}
\question{8}已知 $\int_0^{+\infty} \frac{\sin x^2}{x^p} \, \mathrm{d}x$ 绝对收敛, 则 $p$ 的取值范围为 \_\_\_\_.\par
% p031-b10; source PDF page 31
\par\Needspace{5\baselineskip}
\par\Needspace{6\baselineskip}
\blockanchor{p031-b10}
\majorquestion{二}(10 分) 已知 $D = \{(x, y) \mid x + y \leq 1, x \geq 0, y \geq 0\}$, 且满足\par
\nopagebreak[4]
% p031-b11; source PDF page 31
\blockanchor{p031-b11}
\displayformula{\lim _ {x \rightarrow 1} \left(\frac {\ln x}{x - 1} + \frac {a \sin x + \sin x}{x}\right) = \iint_ {D} \cos \frac {x - y}{x + y} d x d y.}
% p031-b12; source PDF page 31
\blockanchor{p031-b12}
求参数 $a$ 的值.\par
% p031-b13; source PDF page 31
\par\Needspace{5\baselineskip}
\blockanchor{p031-b13}
\majorquestion{三}(10 分) 计算 $\iint_{\Sigma} \frac{x \, \mathrm{d}y \, \mathrm{d}z + y \, \mathrm{d}z \, \mathrm{d}x + z \, \mathrm{d}x \, \mathrm{d}y}{(x^2 + y^2 + z^2)^{\frac{3}{2}}}$, 其中 $\Sigma$ 为 $2x^2 + 2y^2 + z^2 = 4$ 的外侧.\par
% p031-b14; source PDF page 31
\par\Needspace{5\baselineskip}
\blockanchor{p031-b14}
\majorquestion{四}(10 分) 假设 $f(x)$ 在 $(- \infty, + \infty)$ 上一致连续, 证明: 存在非负实数 $a, b$, 使得对任意的 $x \in (-\infty, + \infty)$, 都有 $|f(x)| \leq a|x| + b$.\par
% p031-b15; source PDF page 31
\par\Needspace{5\baselineskip}
\blockanchor{p031-b15}
\majorquestion{五}(10 分) 叙述数列的 Cauchy 收敛准则, 并用 Cauchy 收敛准则证明单调有界数列必收敛.\par
% p031-b16; source PDF page 31
\par\Needspace{5\baselineskip}
\blockanchor{p031-b16}
\majorquestion{六}(10 分) 求函数 $f(x) = x + \frac{1}{1 + x}$ 的单调区间, 凹凸区间, 极值, 拐点和渐进线.\par
% p031-b17; source PDF page 31
\par\Needspace{5\baselineskip}
\blockanchor{p031-b17}
\majorquestion{七}(15 分) 设 $z = z(x, y)$ 由方程 $9x^{2} - 54xy + 90y^{2} - 6yz - 3z^{2} + 18 = 0$ 所确定.\par
% p031-b18; source PDF page 31
\blockanchor{p031-b18}
(1) 求  $\frac{\partial z}{\partial x}, \frac{\partial z}{\partial y}.$\par
% p031-b19; source PDF page 31
\blockanchor{p031-b19}
(2) 求  $z(x,y)$  的极值, 极值点.\par
% p031-b20; source PDF page 31
\par\Needspace{5\baselineskip}
\blockanchor{p031-b20}
\majorquestion{八}(15 分) 设函数 $f(x), g(x)$ 在 $[a, b]$ 上连续, 在 $(a, b)$ 上可导, 且 $f(a) = f(b) = 0$, $g'(x) < 0$. 证明: 至少存在一点 $\xi \in (a, b)$, 使得 $\frac{f'(\xi)}{g'(\xi)} + \frac{\int_a^\xi f(t) \mathrm{d}t}{\int_\xi^b g(t) \mathrm{d}t} = 0$.\par
% p031-b21; source PDF page 31
\par\Needspace{5\baselineskip}
\blockanchor{p031-b21}
\majorquestion{九}(15 分) 设数列 $a_{n}, b_{n}$ 满足 $e^{b_{n}} = e^{a_{n}} - a_{n}, n = 1, 2, \cdots$.\par
% p031-b22; source PDF page 31
\blockanchor{p031-b22}
(1) 若 $a_{n} > 0$, 证明: $b_{n} > 0$.\par
% p031-b23; source PDF page 31
\blockanchor{p031-b23}
(2) 若 $a_{n} > 0$, 正项级数 $\sum_{n=1}^{\infty} a_{n}$ 收敛, 证明: $\sum_{n=1}^{\infty} \frac{b_{n}}{a_{n}}$ 收敛.\par
% p031-b24; source PDF page 31
\par\Needspace{5\baselineskip}
\blockanchor{p031-b24}
\majorquestion{十}(15 分) 讨论下列函数列在给定区间是否一致收敛.\par
% p031-b25; source PDF page 31
\blockanchor{p031-b25}
(1) $f_{n}(x) = \arctan nx, x \in (0, +\infty)$.\par
% p031-b26; source PDF page 31
\blockanchor{p031-b26}
(2) $f_{n}(x) = x\arctan nx, x\in (0, + \infty).$\par
